\fichatitulo{41 · REVISÃO E RISCOS}{Preprint v1 · 2022}{Survey of Hallucination in NLG}
{\small Ji et al. \textit{Survey of Hallucination in Natural Language Generation}. Preprint v1 · 2022. \href{https://arxiv.org/pdf/2202.03629v1}{Fonte consultada}.\par}

\campo{Problema e objetivo}
Como caracterizar e avaliar alucinações em geração de linguagem?

\campo{Principal contribuição · síntese interpretativa}
Organiza o fenômeno em diferentes tarefas e distingue relações com a entrada e com o conhecimento factual.

\campo{Método / natureza da evidência}
Revisão de sumarização, diálogo, perguntas e respostas, geração a partir de dados e tradução.

\campo{Resultado ou argumento principal}
A tipologia separa conteúdo que contradiz a entrada de conteúdo sem suporte nela, oferecendo base para especificar o objeto da avaliação.

\campo{Excertos-chave · seleção literal}
\begin{itemize}
\excerto{1}{Faithfulness in NLG}{Palavras-chave.}
\excerto{2}{Factuality in NLG}{Palavras-chave.}
\end{itemize}

\campo{Limitações e análise crítica}
Análise da ficha: conteúdo não sustentado pela entrada pode ser verdadeiro externamente. Definições operacionais variam por tarefa; esta consulta parcial não audita todas as métricas e técnicas da revisão.

\campo{Aproveitamento na dissertação · proposta}
Fundamentação: definir fidelidade ao contexto e factualidade separadamente, com exemplos de erro parlamentar.

\registro{Fonte consultada nesta preparação: Resumo, organização inicial e exemplos de tipologia. Bibliografia: publicação de 2023; fonte consultada: v1 de 2022. Chave: \texttt{jiEtAl2023hallucinationSurvey}.}
